\documentclass[12pt,a4paper]{article}
\usepackage{ugmath}
\usepackage{float}
\usepackage{placeins}
\usepackage{array}
\usepackage[labelfont=bf,labelsep=space]{caption}
\newcommand{\paren}[1]{\left( #1 \right)}
\newcommand{\alumno}{Ricardo León Martínez}
\newcommand{\materia}{Elementos de Probabilidad y Estadistica}
\newcommand{\profesor}{Claudia Reynoso Alcántara}
\newcommand{\tarea}{Tarea 5}
\newcommand{\fecha}{18/3/2026}

\begin{document}

    \begin{center}
        {\large\textbf{UNIVERSIDAD DE GUANAJUATO}}\\[0.3cm]
        {\normalsize\textbf{DIVISIÓN DE CIENCIAS NATURALES Y EXACTAS}}\\
        {\normalsize\textbf{CAMPUS GUANAJUATO}}\\[1cm]

        {\Large\textbf{\tarea\ (\materia)}}\\[1cm]
    \end{center}

    \noindent
    \textbf{Nombre:} \alumno 
    \hfill 
    \textbf{Fecha:} \fecha 
    \hfill 
    \textbf{Calificación:} \rule{3cm}{0.4pt}

    \vspace{0.3cm}

    \begin{mdframed}[style=mdbluebox,frametitle={Ejercicio 1}]
        En cierta encuesta participaron 1000 personas a las que se les preguntó sobre su ocupación,
        estado civil y nivel educativo. De las personas encuestadas y los datos reportados se sabe que 312
        son profesionistas, 470 están casadas, 525 poseen títulos universitarios, 42 son profesionistas y poseen
        títulos universitarios, 147 están casadas y poseen títulos universitarios, 86 son profesionistas y están
        casadas y 25 son profesionistas casadas que poseen títulos universitarios. Demuestre que los números
        reportados deben ser incorrectos.
    \end{mdframed}

    \begin{mdframed}[style=mdbluebox,frametitle={Ejercicio 2}]
        Un examen contiene 11 preguntas en total a contestar con falso o verdadero. El examen
        resulta aprobado si se responden correctamente al menos 6 de las 11 preguntas. Si contestamos las
        preguntas al azar, ¿cuál es la probabilidad de aprobar el examen?
    \end{mdframed}

    \begin{mdframed}[style=mdbluebox,frametitle={Ejercicio 3}]
        Describa espacios muestrales para cada uno de los siguientes casos.
        \begin{itemize}
            \item[(a)] Calificaciones de una clase con 37 estudiantes en un examen.
            \item[(b)] Medición de la temperatura a medio día en una estación meteorológica.
            \item[(c)] Medición de las velocidades de automóviles pasando por un punto dado.
            \item[(d)] Lanzamientos infinitos de una moneda.
        \end{itemize}
    \end{mdframed}

    \begin{mdframed}[style=mdbluebox,frametitle={Ejercicio 4}]
        Fulgencio y Yesrael acuerdan realizar una competencia de esgrima con duelos consecutivos
        en la que el ganador de la competencia es el primero en salir victorioso en dos duelos consecutivos.
        En cada duelo, la probabilidad de que Fulgencio salga victorioso es $p$ con $0<p<1$, mientras que la
        probabilidad de victoria de Yesrael es $q = 1-p$.
        \begin{itemize}
            \item[(a)] Describa un espacio muestral para este proceso.
            \item[(b)] Calcule la probabilidad de que la competencia termine al cabo de $k$ duelos con $k=2,3,\dots$.
        \end{itemize}
    \end{mdframed}

    \begin{mdframed}[style=mdbluebox,frametitle={Ejercicio 5}]
        Tres jugadores, a los cuales identificamos por Jugador 1, Jugador 2 y Jugador 3, toman
        turnos consecutivos para lanzar una moneda con caras sol y águila. La moneda podría estar desbalanceada,
        por lo que debemos asumir que la probabilidad de que salga águila es $p$ con $0<p<1$, mientras
        que la probabilidad de sol $q = 1-p$. El primer lanzamiento lo realiza el Jugador 1, el segundo el Jugador
        2, el tercero el Jugador 3, el cuarto el Jugador 1, el quinto el Jugador 2, y así sucesivamente. Gana el
        primer jugador que obtenga un sol.
        \begin{itemize}
            \item[(a)] Describa un espacio muestral para este proceso.
            \item[(b)] Describa el evento de que el Jugador 1 gane y calcule su probabilidad.
            \item[(c)] Describa el evento de que el Jugador 2 gane y calcule su probabilidad.
            \item[(d)] Describa el evento de que ni el Jugador 1 ni tampoco el Jugador 2 gana y calcule su probabilidad.
        \end{itemize}
    \end{mdframed}

    \begin{mdframed}[style=mdbluebox,frametitle={Ejercicio 6}]
        Sea $\Omega$ un conjunto no vacío. Recordemos que $\mathcal{F}\subset 2^{\Omega}$ es una
        $\sigma$-álgebra (de subconjuntos) de $\Omega$ si satisface que:
        \begin{itemize}
            \item[(I)] $\Omega \in \mathcal{F}$,
            \item[(II)] $A\in\mathcal{F}$ implica $A^{c}\in\mathcal{F}$,
            \item[(III)] $\{A_k\}_{k=1}^{\infty}\subset \mathcal{F}$ implica $\displaystyle\bigcup_{k=1}^{\infty}A_k \in \mathcal{F}$.
        \end{itemize}
        Más aún, dada una colección $\{S_i\}_{i\in I}\subset 2^{\Omega}$ de subconjuntos de $\Omega$, la $\sigma$-álgebra
        generada por $\{S_i\}_{i\in I}$ es la única $\sigma$-álgebra $\mathcal{F}$ que cumple que:
        \begin{itemize}
            \item[(a)] $\mathcal{F}$ es una $\sigma$-álgebra, es decir, satisface (I), (II) y (III);
            \item[(b)] $S_i\in\mathcal{F}$ para todo $i\in I$;
            \item[(c)] si $\mathcal{F}'$ es una $\sigma$-álgebra tal que $S_i\in\mathcal{F}'$ para todo $i\in I$, entonces $\mathcal{F}\subset \mathcal{F}'$.
        \end{itemize}

        Demuestre la existencia de la $\sigma$-álgebra generada por una colección arbitraria $\{S_i\}_{i\in I}\subset 2^{\Omega}$ de subconjuntos de $\Omega$. Para ello, siga los siguientes pasos:
        \begin{itemize}
        \item[(i)] Note que $2^{\Omega}$ es una $\sigma$-álgebra tal que $S_i\in 2^{\Omega}$ para todo $i\in I$. De este modo, el conjunto
        \[
            \Gamma=\{\mathcal{F}'\subset 2^{\Omega}:\mathcal{F}'\text{ es una }\sigma\text{-álgebra tal que }S_i\in\mathcal{F}'\text{ para todo }i\in I\}
        \]
        es no vacío.
        \item[(ii)] La intersección de todos los elementos de $\Gamma$ pertenece a $\Gamma$, es decir,
        \[
            \bigcap_{\mathcal{F}'\in\Gamma}\mathcal{F}'.
        \]
        \item[(iii)] La $\sigma$-álgebra generada por $\{S_i\}_{i\in I}$ debe ser $\displaystyle\bigcap_{\mathcal{F}'\in\Gamma}\mathcal{F}'$.
        \end{itemize}
        \textbf{Observación.} De manera general, si $\{\mathcal{F}'_j\}_{j\in J}$ es una familia de $\sigma$-álgebras de $\Omega$, entonces
        \[
            \bigcap_{j\in J}\mathcal{F}'_j
        \]
        también es una $\sigma$-álgebra de $\Omega$.
    \end{mdframed}

    \begin{mdframed}[style=mdbluebox,frametitle={Ejercicio 7}]
        Sea $(\Omega,\mathcal{F},P)$ un espacio de probabilidad. Supongamos que
        $\{A_n\}_{n=1}^{\infty}$ y $\{B_n\}_{n=1}^{\infty}$ son dos colecciones numerables de eventos, es decir,
        $A_n,B_n\in\mathcal{F}$ para $n=1,2,\dots$. Demuestre que
        \[
            \lim_{n\to\infty}P(B_n)=1 \quad \text{implica} \quad
            \lim_{n\to\infty}\big(P(A_n)-P(A_n\cap B_n)\big)=0.
        \]

        \textbf{Observación.} Es posible que $\lim_{n\to\infty}P(A_n)$ no exista.
    \end{mdframed}

    \begin{mdframed}[style=mdbluebox,frametitle={Ejercicio 8}]
        Sea $(\Omega,\mathcal{F},P)$ un espacio de probabilidad. Pruebe que se cumple la siguiente desigualdad de Bonferroni:
        \[
        P\left(\bigcup_{k=1}^{n}A_k\right)
        \ge
        \sum_{k=1}^{n}P(A_k)
        -
        \sum_{1\le k<\ell\le n}P(A_k\cap A_\ell).
        \]
        \textbf{Observación/Sugerencia.} Lo anterior es una versión truncada del principio de inclusión-exclusión,
        el cual se prueba por inducción. Recordemos además que para cualesquiera eventos
        $B_1,B_2,\dots,B_n\in\mathcal{F}$ se satisface que
        \[
        P(B_1\cup B_2\cup \cdots \cup B_n)
        \le
        P(B_1)+P(B_2)+\cdots+P(B_n).
        \]
    \end{mdframed}
\end{document}